\section{Introduction}
\label{sec:intro}


Linear Temporal Logic (\ltl) \cite{LTL} is a logic formalism for expressing temporal properties of reactive systems. It extends traditional propositional logic with modal operators, allowing the specification of rules that must hold \emph{through time}, and it is interpreted over traces of discrete events. Variants of \ltl interpreted
over finite traces have been proposed for modeling finite, but length-unbounded, process executions (see Sec. \ref{sec:overview}), making it suitable for finite-horizon problems that often arise in the context of AI and BPM.
Temporal logics have provided one of the most popular tool in Formal Methods for verification of hardware or software system behaviors. In particular, model checking is the problem of verifying whether a model of a system meets a given specification. Despite the importance and the progress made in model checking, it is not always possible to check the model against the specification. This is due to the inherent complexity of the system or the fact that the model is only partially correct or not known. In such cases, one can resort to simpler verification approaches, such as runtime monitoring, where the system's execution is continuously observed on-the-fly and any permanent violation or fulfillment of the specified property is immediately detected. 

Standard symbolic techniques for both model checking and runtime monitoring have achieved great success, but they also suffer for some inherent limitation. Among the others there is the difficulty in dealing with the increasing growth of big data and the lack of flexibility in dealing with highly unpredictable environments, say e.g. evniroments with noisy sensors or probabilistic uncertainties, and adapting the specification when soft violations could happen. Conversely, pure deep learning architectures, such as Recurrent Neural Networks (RNNs) and Transformers, excel at modeling high-dimensional, noisy data distributions, but they likely generate outputs that violate critical safety constraints, as they lack a mechanism to adhere to formal logics. This naturally leads to the use of hybrid approaches, which can possibly overcome some of the mentioned limitations. 

In this paper, we will compare different techniques for runtime monitoring, from purely symbolic methods to the state-of-the-art neuro-symbolic architectures. We start revisiting the neural-symbolic approach to runtime monitoring, originally proposed by \cite{perotti2014neural,perotti2015neural} through the RuleRunner system \cite{perotti2015runtime}. The authors claimed that the resulting neural monitors could exploit GPU parallelism for efficient verification, outperforming purely symbolic mechanisms. Since then, some significant developments have reshaped the landscape. 
First, \ltlf \cite{LTLf} has become the standard formalism for temporal reasoning over finite traces in AI, bringing with it mature tools for compiling formulas into DFAs. Then, differentiable automata representations such as DeepDFA \cite{deepdfa_ecai2024} have emerged, enabling the integration of temporal logic knowledge into neural architectures through gradient-based learning.
\textcolor{red}{Possibly include other NeSy approaches.}
We present a systematic comparison of three paradigms for \ltlf runtime monitoring: (i) purely symbolic DFA-based monitoring, (ii) rule-based neural-symbolic monitoring in the RuleRunner tradition, and (iii) automata-based neural-symbolic monitoring using DeepDFA. 

\noindent
\textcolor{red}{TODO: Mention briefly the experimental results obtained.}